\section{个人理解与开放问题}
\label{sec:discussion}

\subsection{检索必要性是关系属性，而不是问题标签}

综观现有方法，一个容易接受的简化是：简单问题不检索，复杂问题多检索。这一规则在构建路由器时很实用，却不能成为检索必要性的定义。本文更倾向将其理解为四元关系
\[
 \mathcal N(x,\theta,\mathcal D,s_t),
\]
即输入 $x$ 是否需要新证据，取决于当前模型 $\theta$、知识库 $\mathcal D$ 和过程状态 $s_t$。同一个问题交给参数知识不同的模型，检索收益不同；同一个模型连接覆盖范围不同的知识库，最佳路径也不同；同一问题在生成开始时可能需要检索，在获得关键文档后便不再需要。

这一区分带来两个直接推论。首先，只在固定数据集上学习“问题类型—路径”的映射，换模型后可能失效，路由评价应包含跨模型和跨知识库测试。其次，模型“知道事实”与任务“要求可核验引用”是不同需求：即便模型能够回答，高风险场景仍可能为了来源和审计而检索。因此，式~\eqref{eq:objective} 中的质量项应包含任务对可验证性的要求，而不能只以答案字符串是否正确衡量。

\subsection{从模型置信度转向显式证据状态}

现有自适应方法常围绕模型的“不确定性”展开，但本文认为，最终需要控制的并不是模型是否感觉确定，而是答案所依赖的事实是否有足够支持。两者可能分离：模型可以非常自信地说出没有依据的内容，也可以在证据充分时因措辞多样而呈现较高 Token 不确定性。

一种可检验的改进方向是维护显式的证据状态。系统先把回答所需信息拆成若干关键声明或证据槽位，再为每一项记录“缺失、已支持、存在冲突”等状态，同时保存对应文档片段和来源。触发检索时优先选择缺失或冲突项，而不是重新检索整道题；生成后逐项检查输出声明是否能够映射回证据。这里的“已支持”仍可能受到评价模型误判影响，证据槽位也可能拆解不完整，所以它不能消除幻觉，只能让遗漏位置更加可见。

这一设计与单纯扩大上下文有本质区别。目标不是尽可能多地累积 Token，而是减少关键声明的无支持比例，并控制重复证据。它也为模块化组合提供接口：IRCoT 可以发现下一项信息需求，CRAG 可以在证据质量差时修复，Self-RAG 的支持度评价可以参与声明检查。上述组合是本文基于各方法机制提出的研究设想，并非已有工作已经验证的统一系统。

\subsection{停止检索应比较边际证据收益与边际成本}

自适应检索中最被低估的决策是停止。若把第 $t$ 次行动带来的质量改善记为 $\Delta Q_t$，新增检索、生成与风险成本记为 $\Delta C_t$，那么式~\eqref{eq:objective} 的直观含义是：只有当预期收益足以抵消预期成本时，继续行动才合理。实际系统无法精确预知 $\Delta Q_t$，但可以使用可观察代理，例如上一轮是否引入新来源、缺失证据槽位是否减少、冲突是否缓解，以及最近几轮查询是否反复返回相同内容。

据此可以设计分层停止规则。证据已覆盖且无重要冲突时停止；连续检索没有新增有效证据时降低继续搜索的优先级；知识库明显不含答案时进入拒答或澄清，而不是无限改写；达到硬预算时必须停止并披露证据不足。这里的难点是“没有新证据”既可能意味着答案已完整，也可能意味着检索器陷入局部错误，因此停止前可以允许一次有明确目的的替代查询或知识源切换，而非无限制重试。

这一视角也能解释何时迭代方法会退化。如果每轮只重述旧答案、增加重复文档，那么流程虽然形式上在“推理和反思”，证据状态却没有变化；相反，一次准确的查询修正可能已经消除主要缺口。评价时因而应记录每轮新增证据和状态变化，而不只统计总轮数。

\subsection{自我反思需要外部约束，而非被当作事实保证}

Self-RAG、CRAG 及 Agentic RAG 中的反思机制让模型具备评价和修正接口，但“模型评价模型”存在共同失效的可能。若生成器与评价器共享训练数据、提示偏差或对某实体的错误认识，评价器可能认可同样错误的答案；若检索证据本身错误，忠实于证据也只能说明答案与上下文一致，不能证明其与现实一致。

因此，应把几种不同性质的评价分开：相关性检查回答“文档是否对应问题”，忠实性检查回答“输出是否超出文档”，来源检查回答“文档本身是否可信”，跨源一致性检查回答“不同证据是否冲突”。单一反思标记或 LLM 评分很难同时可靠承担这些任务。可探索的方向是保留可追溯引用、对关键声明进行独立抽取验证，并在冲突时展示不确定性。这里的“独立”不仅指再调用一次相同模型，还可以是不同检索来源、规则校验或人工抽查。

\subsection{从方法排名转向可证伪的实验设计}

本文不报告新实验，但根据前述分析，一个能真正比较控制策略的实验矩阵应至少覆盖五组输入：模型熟知且无需检索的问题、长尾或动态知识问题、多跳问题、知识库无答案的问题，以及包含冲突或对抗性噪声的问题。每组分别比较无检索、固定单次检索、固定多轮、基于置信度触发和基于证据状态触发的策略，并统一生成模型、检索器和语料快照。

结果展示不应只有一个平均准确率，而应同时给出质量—检索次数和质量—延迟曲线，以及不必要检索率、漏检率、错误拒答率和无效循环率。特别值得检验的假设包括：复杂度路由器在更换模型后是否仍然有效；低置信度触发是否会在开放写作中造成过度检索；证据槽位能否减少无答案问题上的强行生成；反思模块在生成器与评价器共享错误时是否失效。这样的问题比“再增加一个模块是否提高两个百分点”更能说明自适应机制的边界。

综上，本文对发展方向的判断是：短期最有价值的并非无限扩展 Agent 的自主性，而是让现有闭环变得可观察、可停止和可复核。只有能够解释“为什么查、查到了什么、还缺什么、为什么现在回答”，自适应 RAG 才从一组提示技巧变成真正可靠的证据驱动系统。
